\begin{papercard}{baumann2025}
\cf{Stages} \stagetag{L} (risk of using LLM annotations in analyses; configuration choices; mitigation), \stagetag{I} (human validation samples), \stagetag{R} (pre-registration, reporting all configurations). A large empirical study (arXiv v2, 35 pages plus appendix) that defines ``LLM hacking'': incorrect scientific conclusions caused by researcher choices of model, prompt, decoding parameters and output mapping when LLM annotations feed downstream statistics. It replicates 37 annotation tasks and tests 18 LLMs, giving 13 million annotations and 2,361 hypotheses. The unit of evaluation is the \emph{statistical conclusion}, not annotation accuracy.

\cf{Concept and definitions} ``We call this phenomenon where configuration choices lead to incorrect conclusions LLM hacking'' \pc{baumann2025}{1}{We call this phenomenon where configuration choices lead to incorrect conclusions LLM hacking}. The configuration space is a ``garden of forking paths'' (model, prompt, temperature, token-to-label mapping) \pc{baumann2025}{3}{These choices become a “garden of forking paths”}. It is analogous to p-hacking but acts at the \emph{data-generation} stage rather than at the analysis stage, so the risks accumulate. Formally: LLM hacking occurs when the conclusion under configuration $\phi$ differs from the reference conclusion; \emph{LLM hacking risk} is the probability of an incorrect conclusion when drawing a plausible configuration at random \pc{baumann2025}{6}{The LLM hacking risk quantifies the probability of obtaining false conclusions when randomly selecting from reasonable configurations}; \emph{feasibility} is the existence of at least one configuration producing a desired incorrect conclusion \pc{baumann2025}{6}{We define LLM hacking feasibility as the existence of at least one configuration producing a specific incorrect conclusion}. Error taxonomy: Type I (false positive), Type II (missed effect), Type S (significant effect with the wrong sign), Type M (right sign, distorted magnitude).

\cf{Research questions} RQ1 current practice of LLMs as annotators in computational social science (CSS); RQ2 feasibility of intentional hacking; RQ3 risk of accidental hacking; RQ4 how far mitigation strategies and validation practices reduce the risk.

\cf{Literature review (RQ1)} Systematic search of five databases: 1,592 papers screened to 93 relevant studies (2022--2025). 85 (91.4\,\%) recommend LLM annotation, only 8 advise against; 48 (56.5\,\%) papers recommending LLMs do not validate them; GPT models appear in over 80\,\%; only four mention the risk of false conclusions downstream \pc{baumann2025}{7}{48 (56.5\%) papers disregard model validation (despite recommending the use of LLMs)}.

\cf{Experimental design} \emph{Data}: 37 tasks from 21 public datasets with ground truth (tweets, news, Reddit, survey answers, party manifestos; stance, ideology, frames, hate speech, humour, politeness, \ldots), multiclass tasks decomposed into 250 binary tasks. \emph{Ground truth}: objective labels where available, else expert or crowd labels, removing items on which annotators disagree, giving an average Krippendorff's $\alpha$ of 0.91, so that ``the resulting human annotations'' are treated as noise-free \pc{baumann2025}{8}{we remove all ambiguous datapoints where annotators disagree to ensure high label quality. This filtering process results in a Krippendorff’s alpha of 0.91, on average}. \emph{Hypotheses}: keyword-based splits (texts containing ``economy'' or not) and metadata splits (e.g.\ party, author sex) tested with logistic regression of the (ground-truth or LLM) label on group membership at $\alpha=.05$; 2,361 hypotheses \pc{baumann2025}{10}{We evaluate these hypotheses using logistic regression with binary dependent and independent variables}. \emph{Configurations}: 18 instruction-tuned models (Llama 3 1B--70B, Qwen 2.5/3 1.5B--72B, Gemma 1B--27B, GPT-4o variants), 199 prompts (5--7 per task; 36\,\% paraphrases, 89\,\% zero-shot), temperature 0, max 20 tokens, regex mapping to labels \pc{baumann2025}{11}{We evaluate 18 models from 4 model providers representing the current state of the art}. \emph{Metric}: weighted F1 for annotation quality, risk measures by error type averaged over tasks and configurations. \emph{Baselines}: random conclusions, random labels, random errors with controlled F1 \pc{baumann2025}{12}{Baseline 3 (random errors): Generates annotations with controlled F1 scores by flipping random labels to incorrect classes}.

\cf{Mitigation techniques (RQ4)} 21 variants combining sampling of items for human labelling (random, low LLM confidence, active learning with XGBoost-predicted error), use of human labels (human only; human plus LLM; human plus LLM with corrected estimator: Design-based Supervised Learning, DSL, and Confidence-Driven Inference, CDI), and model selection (random, GPT-4o, best on the validation sample).

\cf{Results: intentional hacking} Using at least one of the tested model--prompt pairs, a false positive can be produced for 94.4\,\% of null hypotheses, a true effect hidden in 98.1\,\%, and a significant effect reversed (Type S) in 68\,\% of cases with true effects; even restricting to the seven best models the rates remain about 70\,\% \pc{baumann2025}{16}{Even when restricting analyses to top-performing models – which a reviewer might consider methodologically sound – the intentional hacking feasibility remains alarmingly high}. Constraining the configuration space or requiring replicability across independent configurations reduces feasibility, hence ``Pre-registration of all LLM configuration choices, combined with reporting results from all pre-registered configurations'' \pc{baumann2025}{17}{Pre-registration of all LLM configuration choices, combined with reporting results from all pre-registered configurations, thus makes intentional LLM hacking more difficult}.

\cf{Results: accidental hacking} Risk of an incorrect conclusion: about 31\,\% for the best 70B models up to 50\,\% for 1B models; by task from 5\,\% (humour) to above 65\,\% (ideology, factuality, frames); Type II errors dominate (31--59\,\%), Type S 4--16\,\%, Type M 41--77\,\% \pc{baumann2025}{18}{empirical LLM hacking risk ranges from 31\% for the best-performing 70B parameter models to 50\% for 1B parameter models}. Random labels have very low Type I risk (7.8\,\%) but Type II risk 92.9\,\% \pc{baumann2025}{18}{Random labels have extremely low Type I risk (7.8\%)}. Risk factors (regression on all hypothesis--model--prompt cells): closeness of the effect to the significance threshold is the strongest predictor (error rates near 70\,\% at $p\approx.05$), task characteristics explain about 21\,\% of the explained variance, model performance 8\,\%, prompt choices under 1\,\%; \textbf{no correlation} between human inter-annotator agreement and LLM hacking risk \pc{baumann2025}{4}{Surprisingly, we find no correlation between human inter-annotator agreement and LLM hacking risk}.

\cf{Results: mitigation} ``100 human annotations outperform 100K LLM annotations for controlling false positives'': human labels alone give Type I error around 10\,\% with 100 labels, vs.\ 30--40\,\% for uncorrected human+LLM mixes; DSL/CDI restore nominal Type I error but raise Type II error by up to 60 percentage points; even 1,000 human labels leave overall risk near 20\,\%; testing several models and choosing the best on a sample yields only minor gains \pc{baumann2025}{5}{Regression estimator correction techniques restore valid inference but trade Type I for Type II errors}.

\cf{Recommendations (Table 5, summarised)} Task-unspecific: largest models (70B+ about 27\,\% risk vs.\ 47\,\% for 1B), extreme caution near $p=.05$, few-shot over zero-shot, detailed task descriptions. Human-sample-enabled: collect as many expert labels as possible; random human samples give lowest Type I error; correction methods trade error types; avoid low-accuracy annotations; do not use human agreement to decide whether to automate \pc{baumann2025}{34}{Do not rely on human annotator agreement to decide whether annotations should be automated or not}. Transparency: report all models, versions, prompts and parameters tested, selection criteria, release LLM and human labels with code; pre-register model/prompt/parameter selection criteria, hypotheses, tests and sensitivity analyses \pc{baumann2025}{34}{Report all models, versions, prompts, and parameters tested}. ``Document all tested model-prompt combinations, not just the final choice'' \pc{baumann2025}{18}{Document all tested model-prompt combinations, not just the final choice}. Future: multiverse analyses that report coefficient distributions across reasonable configurations \pc{baumann2025}{32}{reporting coefficient distributions across all reasonable LLM configurations rather than single point estimates}; human-in-the-loop designs.

\cf{Limits (authors and card)} Assumes noise-free ground truth with the disagreement items removed (real human labels contain error and bias) \pc{baumann2025}{31}{Our LLM hacking risk quantification assumes noise-free ground truth annotations and effective significance thresholds}; risk depends on the plausible configuration space (e.g.\ misspecified prompts could inflate it); outcome defined as a significance decision of a logistic regression at $p<.05$ on binary labels, which differs from other analyses (descriptive prevalence, effect sizes, qualitative coding); mostly English social-media/news/manifesto text; a one-pass, temperature-0, short-output design; models of 2024--25; the 94 \% / 98 \% feasibility figures are best-case for a manipulator across many configurations and are not typical error rates; synthetic keyword and metadata hypotheses rather than hypotheses of real studies.

\cf{Use for this thesis} (1) The strongest argument for how to \emph{report} stage L: pre-register (in the Decision Log) the model, version, prompt, temperature, mapping and selection rule before looking at the test labels; report every configuration tried; show a small multiverse (several prompt paraphrases, two or more models) rather than one run. (2) Evidence that high human agreement is not a licence to automate, and that near-threshold findings are fragile; for the thesis this means LLM labels may be used for screening and prevalence description only after validation, and no inferential claim should rest on LLM labels alone. (3) Quantitative support for keeping a human-coded validation sample (100--1,000 items) rather than adding more LLM runs. (4) Avoid transferring the numbers (31--50\,\%) to the thesis: they concern significance testing on social-media tasks, not coding of privacy features.
\end{papercard}
